\section{Measurable maximal segments and cutoff-uniform source bounds}
\label{sec:v85-segments}
All surfaces in this article are counted with their incident physical
component. The following parametrization makes this convention, and the
absence of a word-count factor, explicit. Fix $m,R,n,k$ throughout its
construction. Let $\Omega$ be the regular exact event, with the critical
points of $F$ removed, and let $\tau^\pm$ be the two lifetimes of $X$.
Write $\bar\nu=\nu_R^*|_\Omega$.

\begin{theorem}[A Borel parametrization of maximal regular segments]
\label{thm:v85-segment-map}
There are countably many Borel sets $E_j$ in lifted angular intervals,
Borel functions $a_j<b_j$, and a Borel bijection
\begin{equation}\label{eq:v85-segment-map}
 \Psi:\coprod_j\{(\theta,u):\theta\in E_j,\ a_j(\theta)<u<b_j(\theta)\}
                  \longrightarrow\Omega
\end{equation}
with Borel inverse, such that each $\theta\in E_j$ labels one entire
maximal segment, $F(\Psi_j(\theta,u))=u$, and
$\vartheta(\Psi_j(\theta,u))=\theta$ modulo $2\pi$. Its physical
angular density is
\[
 A_j(u,\theta)=\frac{\omega(\Psi_j(\theta,u))}
                         {4\pi c_*\chi(\Psi_j(\theta,u))}.
\]
For every nonnegative Borel function $f$,
\begin{equation}\label{eq:v85-segment-area}
 \int_\Omega f\,d\bar\nu
   =\sum_j\int_{E_j}\int_{a_j(\theta)}^{b_j(\theta)}
           f(\Psi_j(\theta,u))A_j(u,\theta)\,du\,d\theta.
\end{equation}
The same identity holds with any Borel selector of the segment and its
interior point. For almost every $t$, disintegration gives the
corresponding identity on $\mu_t$, with $u=t$. Both identities are
identities of measures, not estimates with an unspecified multiplicity.

Restrict to segments with a nongrazing forward endpoint. The endpoint
map, retaining the incident one-sided component, is Borel and injective
outside the transverse null sets of the first-exit construction. Its
outgoing measure in these coordinates is
\begin{equation}\label{eq:v85-segment-face}
                    d\zeta_+=A_j(b_j-,\theta)\,d\theta.
\end{equation}
Each such face and an intermediate roof determine at most one physical
interior state. In particular, summing a variation integral over these
segments gives exactly the selected physical volume integral in
\eqref{eq:v85-segment-area}.
\end{theorem}
\begin{proof}
Enumerate the rational numbers as $q_1,q_2,\ldots$. The lifetime
functions of a smooth flow on an open set are lower semicontinuous;
this also follows by covering each compact surviving segment with
finitely many ordinary flow boxes. They are therefore Borel. Since
$XF=1$ and $F$ is bounded at fixed count, the roof interval of every
maximal segment is a bounded, nonempty open interval.

On the smooth transverse slice $\Sigma_i=\Omega\cap\{F=q_i\}$,
keep the points whose maximal interval contains none of
$q_1,\ldots,q_{i-1}$. These are Borel conditions on the lifetimes.
Every maximal segment has exactly one representative in this union:
use the first rational in its roof interval and strict monotonicity of
$F$. On each slice, $\vartheta$ is a local coordinate, since
$\det D_q(F,\vartheta)=\chi>0$. Cover the slices by countably many
angular charts, and disjointize their representatives by the first
containing chart. Denote the resulting Borel angular subsets by $E_j$,
and the associated representative by $\xi_j(\theta)$, at roof $q_{i(j)}$.
Put
\[
 a_j=q_{i(j)}-\tau^-(\xi_j),\qquad
 b_j=q_{i(j)}+\tau^+(\xi_j),\qquad
 \Psi_j(\theta,u)=\Phi_{u-q_{i(j)}}\xi_j(\theta).
\]
The flow is continuous on its open lifetime domain. This proves Borel
measurability. The inverse is explicit: for $x$, select the first
rational in its maximal roof interval, flow to that roof, and select
the first angular chart containing the representative. All these
selections are Borel. If two parameter points have the same image,
uniqueness of the flow gives the same maximal segment and representative;
then $F=u$ and the angular chart give identical parameters. This proves
\eqref{eq:v85-segment-map} without a quotient measurability assumption.
The Borel seams assign representatives only; they are not new exits.

In a local angular flow box, change of variables from endpoint angles
has density $\omega/(4\pi c_*\chi)$. The same formula at another roof
follows by physical coarea covariance; the transverse angle is unchanged.
Apply this local formula on countably many Borel-disjoint restrictions
of the parametrization and use Tonelli. It gives
\eqref{eq:v85-segment-area} for every Borel selector. Disintegration over
$u$ and a countable determining algebra give equality of fiber measures
outside one null set in $t$, then equality for all nonnegative Borel
functions. Additional external variables use the product graph and
Fubini, as in Proposition~\ref{prop:v83-common-kernel}.

The endpoint limits exist by the all-contact motion estimate. They are
Borel as limits along the interior roofs
$b_j-\min\{1/N,(b_j-a_j)/2\}$. A nongrazing forward endpoint has
positive limiting gradient and a one-sided smooth stationary-word
extension. The flow-box contraction of the physical area form gives
\eqref{eq:v85-segment-face}. Two coincident geometric boundary points
on different incident components are different faces by definition.
On one incident component, uniqueness of the backward regular flow
makes its interior segment unique. Tangent and lower-dimensional
strata have the transverse-null convention of
Proposition~\ref{prop:v83-first-exit-kernel}; removing their flow
saturations removes no area in \eqref{eq:v85-segment-area}. This proves
the last assertions. No continuation past a true physical boundary is
used, and a backward critical endpoint remains an endpoint of its
segment, not a chart seam.
\end{proof}

\subsection{The lower cutoff and its finite-count remainder}
We record the quantitative step behind the cutoff-independent constant
used in boundary transfer. It concerns an actual positive volume source;
it is not a surface insertion theorem. Write $z=|g|$, $\alpha=129/128$,
and
\[
 u^b(t)\,dt=F_\#(I_m\one_{\{z<b\}}\bar\nu),\quad
 v^{c,b}(t)\,dt=F_\#(I_mz^{-1}\one_{\{c\le z<b\}}\bar\nu),\quad
 v^c(t)\,dt=F_\#(I_mz^{-1}\one_{\{z\ge c\}}\bar\nu).
\]
The source-pinned critical-volume estimate, including its remainder, is
\begin{equation}\label{eq:v85-critical-input}
 m^2\int_Ju^b\le Cb^\alpha(|J|+B_*^{-1})
 +C_{B_*}(1+B_*|J|)m^2b^{1/2}\rho_{B_*}^{m/4}.
\end{equation}
It is Supplement equation~\eqref{supp-v81archive:eq:v81-critical-finite}.
Its dependence on the joint spectral theorem is listed in the dependency
map; it is restated here, not independently recertified.

\begin{proposition}[Uniform leading constant at every positive cutoff]
\label{prop:v85-cutoff}
For $0<c<b<b_1$, the complete exact source satisfies
\begin{align}
 m^2\int_Jv^{c,b}
 \le{}&C\left(1+\frac1{\alpha-1}\right)b^{\alpha-1}
                                      (|J|+B_*^{-1})\notag\\
 &+2C_{B_*}(1+B_*|J|)m^2c^{-1/2}\rho_{B_*}^{m/4}.
                                             \label{eq:v85-cutoff-finite}
\end{align}
Consequently, for every fixed $c,L>0$,
\begin{equation}\label{eq:v85-cutoff-uniform}
 \limsup_m\sup_{R,n,k}\sup_{|J|=L}m^2\int_Jv^c\le C_0 L,
 \qquad C_0\text{ independent of }c\text{ and }L.
\end{equation}
No assertion at $c=0$, or at a collision-dependent $c$, is made.
\end{proposition}
\begin{proof}
Use the positive majorant
\[
 z^{-1}\one_{\{c\le z<b\}}
 \le b^{-1}\one_{\{z<b\}}+\int_c^b q^{-2}\one_{\{z<q\}}\,dq.
\]
Insert \eqref{eq:v85-critical-input} before any limit. The leading
coefficient is at most
$b^{\alpha-1}+\int_c^b q^{\alpha-2}\,dq
\le(1+(\alpha-1)^{-1})b^{\alpha-1}$.
The error coefficient is
$b^{-1/2}+\int_c^b q^{-3/2}\,dq
=2c^{-1/2}-b^{-1/2}\le2c^{-1/2}$.
This proves \eqref{eq:v85-cutoff-finite}. Split $v^c$ at the fixed
threshold $b_1/2$. Above it, use the complete depth-weighted source
estimate of Supplement Corollary~\ref{supp-v80archive:cor:v80-depth-source}.
Below it, apply \eqref{eq:v85-cutoff-finite}, first fixing $c,L,B_*$,
then taking the collision limsup, and then enlarging $B_*$. If
$c\ge b_1/2$, positivity gives the same conclusion. The surviving
leading constant contains only fixed physical thresholds, not $c$.
The displayed finite-count error explains why removal of $c$ before
the collision limit is a separate question.
\end{proof}

\begin{corollary}[The audited paired-transfer estimate]
\label{cor:v85-audited-transfer}
On any Borel subfamily of short segments with $b-a\le2r$ and
$z(b-)^2\ge8\Lambda_\circ r$, let $S_r^+$ and $S_r^-$ use the
same exit-window kernel and the outgoing and incoming angular traces,
respectively. Then
\[
 |S_r^+(t)-S_r^-(t)|\le C\int_{t-2r}^{t+2r}v^{c_r}(u)\,du,
 \qquad c_r=2\sqrt{\Lambda_\circ r},
\]
and their ordered global essential-height difference is at most
$4CC_0r$, with $C_0$ independent of $c_r>0$. A grazing incoming
trace is zero, so its entire outgoing contribution obeys this bound.
\end{corollary}
\begin{proof}
Radial growth keeps the full segment above $c_r$. Differentiate the
physical angular density, retaining both factors:
$|\partial_u A|\le C A I_m/z$, as proved in Supplement
Lemma~\ref{supp-v84archive:lem:v84-angular-variation}. At a grazing
incoming end $\omega\to0$ while $\chi\ge\lambda_\circ$; at any
nongrazing end the one-sided regular extension gives its trace.
Integrate the absolute variation between these traces. When its kernel
is nonzero the full selected segment lies in $(t-2r,t+2r)$.
Theorem~\ref{thm:v85-segment-map} identifies the complete sum with a
selected submeasure of $v^{c_r}(u)\,du$, with multiplicity one.
Positive approximation proves the assertion for all selected sheets.
Apply \eqref{eq:v85-cutoff-uniform} with $L=4r$. This supplies the
named map, exact measure, and cutoff uniformity in the same argument.
\end{proof}
